\subsection{\texorpdfstring{The semi-norms \(\mathbf{N}_p\)}{The semi-norms N sub p}}

In all that follows, F denotes a complete normed vector space (that is, a Banach space) over the field R or the field C; the norm of an element \(z \in F\) will be denoted \(|z|\) . Given a mapping f of a set A into F, we write \(|f|\) for the mapping \(x \mapsto |\mathbf{f}(x)|\) of A into \(R_{+}\) (one must take care to note that \(|f|\) is a numerical function, and not a number).

DEFINITION 1. --- Let X be a locally compact space, \(\mu\) a measure on X. For every mapping f of X into a Banach space F, and for every number p such that \(1 \leqslant p < +\infty\) , we denote by \(\mathrm{N}_p(\mathbf{f}, \mu)\) , or simply \(\mathrm{N}_p(\mathbf{f})\) , the positive number \((\int^* |\mathbf{f}|^p d|\mu|)^{1/p}\) .

Note that the number \(N_{p}(\mathbf{f})\) may be equal to \(+\infty\) .

PROPOSITION 2. --- If f and g are two mappings of X into F, and \(\alpha\) is any scalar \(\neq 0\) , then, for \(1 \leqslant p < +\infty\) ,

\begin{enumerate}
\def\labelenumi{(\arabic{enumi})}
\setcounter{enumi}{2}
\tightlist
\item
\end{enumerate}

\[
\mathrm{N} _ {p} (\alpha \mathbf {f}) = | \alpha | \mathrm{N} _ {p} (\mathbf {f})\tag{4}
\]

\[
\mathrm{N} _ {p} (\mathbf {f} + \mathbf {g}) \leqslant \mathrm{N} _ {p} (\mathbf {f}) + \mathrm{N} _ {p} (\mathbf {g}).
\]

For, the relation (3) follows at once from Def. 1 and the fact that \(|\mu|^*\) is positively homogeneous; on the other hand, since \(|\mathbf{f} + \mathbf{g}| \leqslant |\mathbf{f}| + |\mathbf{g}|\) , the inequality (4) follows from Minkowski's inequality (1) and the fact that \(|\mu|^*\) is increasing.

We extend Def. 1 to the case of numerical functions, finite or not, defined on X, by again setting

\[
\mathrm{N} _ {p} (f) = \left(\int^ {*} | f | ^ {p} d | \mu |\right) ^ {1 / p}
\]

for such a function \(f\) . One sees immediately that the relations (3) and (4) also hold for these functions when \(f + g\) is defined on \(X\) and \(\alpha \neq 0\) . Moreover:

THEOREM 1 (countable convexity theorem). --- Let \((f_n)\) be a sequence of functions \(\geqslant 0\) (finite or not) defined on X. For \(1 \leqslant p < +\infty\) ,

\[
\mathrm{N} _ {p} \left(\sum_ {n = 1} ^ {\infty} f _ {n}\right) \leqslant \sum_ {n = 1} ^ {\infty} \mathrm{N} _ {p} (f _ {n}).\tag{5}
\]

Set \(f = \sum_{n=1}^{\infty} f_n\) ; \(f\) is the upper envelope of the increasing sequence of functions \(g_n = \sum_{k=1}^{n} f_k\) ; the definition of \(N sub p(f)\) , and Th. 3 of §1, No.~3, show that \(\mathrm{N}_{p}(f)=\sup_{n}\mathrm{N}_{p}(g_{n})\) . But \(\mathrm{N}_{p}(g_{n})\leqslant\sum_{k=1}^{n}\mathrm{N}_{p}(f_{k})\) by Prop. 1, whence the inequality (5).

PROPOSITION 3. --- If \(\mathbf{f}\) and \(\mathbf{g}\) are two equivalent mappings of X into a Banach space F, then \(\mathrm{N}_p(\mathbf{f} - \mathbf{g}) = 0\) for \(1 \leqslant p < +\infty\) ; conversely, if \(\mathrm{N}_p(\mathbf{f} - \mathbf{g}) = 0\) for a value of \(p \geqslant 1\) then \(\mathbf{f}\) and \(\mathbf{g}\) are equivalent. The proposition follows at once from Th. 1 of §2, No.~3.

If \(\mathbf{f}\) and \(\mathbf{g}\) are two equivalent mappings of \(X\) into \(F\) , then \(\mathrm{N}_p(\mathbf{f}) = \mathrm{N}_p(\mathbf{g})\) for all \(p \geqslant 1\) ( \(\S 2\) , No.~3, Prop. 6); thus, \(\mathrm{N}_p(\mathbf{f})\) depends only on the class \(\widetilde{\mathbf{f}}\) of \(\mathbf{f}\) , and one sets, by definition, \(\mathrm{N}_p(\widetilde{\mathbf{f}}) = \mathrm{N}_p(\mathbf{f})\) . Since the classes of mappings of \(X\) into \(F\) form a vector space ( \(\S 2\) , No.~4), the relations (3) and (4) may also be written

\begin{enumerate}
\def\labelenumi{(\arabic{enumi})}
\setcounter{enumi}{5}
\tightlist
\item
\end{enumerate}

\[
\mathrm{N} _ {p} (\alpha \widetilde {\mathbf {f}}) = | \alpha | \mathrm{N} _ {p} (\widetilde {\mathbf {f}})\tag{7}
\]

\[
\mathrm{N} _ {p} (\widetilde {\mathbf {f}} + \widetilde {\mathbf {g}}) \leqslant \mathrm{N} _ {p} (\widetilde {\mathbf {f}}) + \mathrm{N} _ {p} (\widetilde {\mathbf {g}}).
\]

One defines similarly \(\mathrm{N}_{p}(\tilde{\mathbf{f}})\) for every class of equivalent numerical functions (finite or not).

One can then define \(\mathrm{N}_{p}(\mathbf{f})\) for a function with values in F (resp. in \(\overline{R}\) ) defined almost everywhere in X, by setting \(\mathrm{N}_{p}(\mathbf{f})=\mathrm{N}_{p}(\widetilde{\mathbf{f}})\) ; it is then clear that the relations (3) and (4) again hold (assuming \(\alpha\neq0\) and \(f+g\) defined almost everywhere, in the case of numerical functions, finite or not).

If \(0 < p < 1\) , one again sets \(\mathrm{N}_p(\mathbf{f}) = (\int^* |\mathbf{f}|^p d|\mu|)^{1/p}\) , but the inequalities (4) and (5) are no longer valid (cf.~Ch. I, Exer. 6 and Ch. IV, §6, Exer. 13).

